% Part: many-valued-logic
% Chapter: syntax-and-semantics
% Section: formulas

\documentclass[../../../include/open-logic-section]{subfiles}

\begin{document}

\olfileid{mvl}{syn}{fml}

\olsection{公式}

\begin{defn}[公式]
\ollabel{defn:formulas}
命题语言~$\Lang L$ 的\emph{公式}集合~$\Frm[L]$ 归纳定义如下：
\begin{enumerate}
\item 每个命题变元~$\Obj p_i$ 都是原子公式。
\item $\Lang L$ 的每个 $0$ 元联结词（命题常项）都是原子公式。
\item 如果 $\star$ 是 $\Lang L$ 的 $n$ 元联结词，且 $!A_1$,
\dots, $!A_n$ 是公式，则 $\star(!A_1, \dots, !A_n)$ 是公式。
\tagitem{limitClause}{除此之外皆非公式。}{}
\end{enumerate}
如果 $\star$ 是一元联结词，则 $\star(!A_1)$ 通常简写为 $\star !A_1$。如果 $\star$ 是二元联结词，$\star(!A_1,!A_2)$
通常简写为 $(!A_1 \star !A_2)$。
\end{defn}

按照惯例，我们通常会省略最外层的括号。

\begin{ex}
  在标准语言~$\Lang{L_0}$ 中，$\Obj p_1 \lif (\Obj p_1
  \land \lnot \Obj p_2)$ 是公式。在积逻辑的语言中，则会写作 $\Obj p_1 \lif (\Obj p_1 \odot
  \lnot \Obj p_2)$。如果在语言中加入一元联结词 $\triangle$，则还会有诸如 $\triangle (\Obj p_1
  \land \Obj p_2) \lif (\triangle \Obj p_1 \land \triangle \Obj p_2)$ 这样的公式。
\end{ex}

\end{document}